% !TeX root = ../../Book3.tex
\section{双曲型、抛物型与椭圆型}
\label{b3:sec:1203}

共形类型指万有覆盖的模型，不是底曲面是否有界、是否紧或是否有孔的简单同义词。特别是紧曲面也可能有圆盘万有覆盖。

\subsection{共形类型}
\label{b3:subsec:120301}

\begin{definition}\label{b3:def:12-conformal-types}
万有覆盖分别为球面、平面、圆盘的曲面称为\term[tuoyuanxing-qumian]{椭圆型曲面}[elliptic Riemann surface]、\term[paowuxing-qumian]{抛物型曲面}[parabolic Riemann surface]、\term[shuangquxing-qumian]{双曲型曲面}[hyperbolic Riemann surface]。这三个词在本卷始终按覆盖模型解释。
\end{definition}
势论中的“有无 Green 函数”也常称为双曲与抛物分类，但对一般非单连通曲面，它与这里的覆盖分类不宜无条件等同。紧的高亏格曲面具有圆盘覆盖，却没有单个对数极点的全局 Green 函数；紧性使 Laplace 方程需要额外的零平均补偿。

\ProseParagraphBreak
本书用“覆盖双曲型”或明确提到圆盘覆盖消除这个语境差别。不要只由一个词的英文相同就把不同分类中的性质互相搬用。

\ProseParagraphBreak
曲面类型说的是万有覆盖，并不要求曲面本身等于那个模型。有限环域、去心盘及紧亏格至少2曲面都为圆盘型，却分别有两条普通边界、一个穿孔端或没有边界的紧拓扑。平面型也包含去心平面和紧环面，而非只包含平面。术语若用在一般曲面，必须把覆盖定义保留；对于单连通曲面才可直接把它与模型等同。

\ProseParagraphBreak
尤其不能把本章的抛物型直接读成“没有 Green 函数”的全部一般曲面分类。对单连通非紧模型两种观点相互对应，但对一般商曲面，势论是否存在 Green 函数还与商的全局结构有关。本章只按覆盖定义命名，后面的势论会按其自己的存在条件陈述。相同形容词在不同分类框架中的定义范围必须看清。

\subsection{万有覆盖}
\label{b3:subsec:120302}

平面甲板群对应一个离散加法子群 \(\Lambda\subset\mathbb C\)。局部不交性给原点附近没有非零群元。若没有非零元，商为平面；若群位于一条实直线上，最短正方向元生成它，商为去心平面；若存在两个实线性无关元，取其格的基本平行四边形，离散性使其中只有有限候选，再缩减可选出一个格基，群为 \(\mathbb Z\omega_1+\mathbb Z\omega_2\)。其商为复环面。

\ProseParagraphBreak
因此抛物型曲面恰有平面、去心平面、复环面这三类形式；复环面仍有不同模参数，绝非只有一个共形类。球面型只有球面。其他曲面的覆盖模型只能为圆盘。这个分类由自由甲板变换与离散平移群得到，不需要假定底曲面有一个平面单值坐标。

\subsection{例子}
\label{b3:subsec:120303}

标准有限圆环的万有覆盖是带，带共形等价于圆盘，所以有限圆环双曲型。去心圆盘覆盖为半平面，也双曲型；去心平面覆盖为全平面，因此抛物型。三者基本群都同构 \(\mathbb Z\)，故基本群的抽象同构类不能区分它们的共形类型。

\ProseParagraphBreak
复环面 \(\mathbb C/\Lambda\) 紧而抛物型；球面紧而椭圆型；亏格至少2的紧曲面双曲型。最后一句可由其基本群非交换排除平移商，也由非球面拓扑排除球面。底曲面紧性只决定商是否紧，并不决定万有覆盖是否紧。

\subsection{穿孔与紧曲面}
\label{b3:subsec:120304}

从球面删去一点得平面，删去两点得去心平面，删去至少三点则为双曲型平面区域。这展示了穿孔数量改变覆盖模型的临界现象。对复环面删去一点，得到非紧且基本群非交换的曲面，也双曲型。

\ProseParagraphBreak
“椭圆型”与“椭圆函数”不是同一个分类。椭圆函数生活在平移格商复环面上，其万有覆盖是平面，按本卷定义为抛物型。名称的历史来源不同，不能据字面把环面放入球面型。类似地，“双曲型”也不表示曲面作为欧氏图形有双曲线边界。

\ProseParagraphBreak
穿孔可能改变类型。球面删去一点是平面，删去两点是 \(\mathbb C^*\)，删去三点已为圆盘型。三点的位置可用 Möbius 变换归一化为 \(0,1,\infty\)，得到 \(\mathbb C\setminus\{0,1\}\)。基本群含两条独立绕圈，并非平面型商的秩0或秩1平移群；同时区域不紧，不能是秩2平移商的环面。这与一般平面双曲分类的证明相符。

\ProseParagraphBreak
环面删去一点也变成圆盘型。若仍是平面型，上述分类只允许平面、去心平面或完整环面；前三者的拓扑分别为无孔、一个基本圈与紧面，不等于去掉一点的环面。这里使用曲面的通常拓扑分类或基本群计算。没有把局部删去一个圆盘的图误当作整个曲面；穿孔附近看起来都像去心盘，但全局剩余的柄仍影响覆盖类型。

\ProseParagraphBreak
亏格至少2的紧曲面不能是球面型，也不能是平面型：前者只有球面，后者的唯一紧商是环面。因而由统一化它必为圆盘型。亏格1的紧曲面反向确为环面，需要曲面的全局函数或1形式理论；附录A从 Riemann--Roch 给出这一联系。若只知道一个给定商 \(\mathbb C/\Lambda\)，其环面结构直接可见，不需要反向分类。区分“所构造模型属于某类型”与“所有这种拓扑的复结构都属于某类型”，可避免在例子中遗漏真正的存在论证。
